\chapter{Variational Autoencoders}

\section{Latent Representations}

Variational autoencoders (VAEs) are generative models that combine latent variable modeling with neural networks.

\medskip

\noindent
\textbf{Core idea.}  
Learn a low-dimensional latent representation $z$ that captures the underlying structure of data $x$.

\medskip

\noindent
\textbf{Generative model.}
\[
z \sim p(z), \quad x \sim p_\theta(x \mid z),
\]
where $p(z)$ is typically chosen as a standard Gaussian.

\medskip

\noindent
\textbf{Latent space.}
\begin{itemize}
    \item Continuous, low-dimensional representation
    \item Encodes semantic structure of data
\end{itemize}

\medskip

\noindent
\textbf{Interpretation.}
\begin{itemize}
    \item Nearby points in latent space correspond to similar data
    \item Interpolation in latent space produces meaningful variations
\end{itemize}

\medskip

\noindent
\textbf{Goal.}  
Learn both:
\begin{itemize}
    \item a generative model $p_\theta(x \mid z)$,
    \item an inference model $q_\phi(z \mid x)$.
\end{itemize}

\medskip

\noindent
\textbf{Insight.}  
Latent representations provide a structured way to model complex data distributions.

\section{Evidence Lower Bound (ELBO)}

VAEs are trained by maximizing the evidence lower bound introduced in the previous chapter.

\medskip

\noindent
\textbf{ELBO objective.}
\[
\mathcal{L}(\theta, \phi)
=
\mathbb{E}_{q_\phi(z \mid x)}[\log p_\theta(x \mid z)]
-
\mathrm{KL}(q_\phi(z \mid x) \,\|\, p(z)).
\]

\medskip

\noindent
\textbf{Interpretation.}
\begin{itemize}
    \item First term: reconstruction accuracy
    \item Second term: regularization toward the prior
\end{itemize}

\medskip

\noindent
\textbf{Tradeoff.}
\begin{itemize}
    \item Good reconstruction vs structured latent space
\end{itemize}

\medskip

\noindent
\textbf{Reparameterization trick.}  
To optimize the expectation, we express $z$ as:
\[
z = \mu_\phi(x) + \sigma_\phi(x) \odot \varepsilon, \quad \varepsilon \sim \mathcal{N}(0, I).
\]

\medskip

\noindent
\textbf{Benefit.}
\begin{itemize}
    \item Enables gradient-based optimization
    \item Reduces variance of gradient estimates
\end{itemize}

\medskip

\noindent
\textbf{Insight.}  
The ELBO balances data reconstruction and latent structure, forming the core of VAE training.

\section{Encoder--Decoder Architectures}

VAEs are implemented using neural networks in an encoder--decoder framework.

\medskip

\noindent
\textbf{Encoder.}
\[
q_\phi(z \mid x)
\]
maps input data to latent parameters $(\mu_\phi(x), \sigma_\phi(x))$.

\medskip

\noindent
\textbf{Decoder.}
\[
p_\theta(x \mid z)
\]
maps latent variables back to the data space.

\medskip

\noindent
\textbf{Architecture.}
\begin{itemize}
    \item Encoder network: compresses data into latent representation
    \item Decoder network: reconstructs data from latent code
\end{itemize}

\medskip

\noindent
\textbf{Training pipeline.}
\begin{enumerate}
    \item Encode $x$ to $(\mu, \sigma)$
    \item Sample $z$ using reparameterization
    \item Decode $z$ to reconstruct $x$
    \item Compute ELBO and update parameters
\end{enumerate}

\medskip

\noindent
\textbf{Sampling.}  
New data is generated by:
\begin{itemize}
    \item Sampling $z \sim p(z)$
    \item Computing $x \sim p_\theta(x \mid z)$
\end{itemize}

\medskip

\noindent
\textbf{Insight.}  
The encoder--decoder structure integrates inference and generation into a unified framework.

\section{Practical Considerations}

Implementing VAEs involves several important design choices.

\medskip

\noindent
\textbf{Choice of likelihood.}
\begin{itemize}
    \item Gaussian likelihood for continuous data
    \item Bernoulli likelihood for binary data
\end{itemize}

\medskip

\noindent
\textbf{KL divergence term.}  
Often has a closed form when both $q_\phi(z \mid x)$ and $p(z)$ are Gaussian.

\medskip

\noindent
\textbf{Posterior collapse.}
\begin{itemize}
    \item The model ignores latent variables
    \item $q_\phi(z \mid x)$ becomes close to $p(z)$
\end{itemize}

\medskip

\noindent
\textbf{Mitigation strategies.}
\begin{itemize}
    \item KL annealing
    \item More expressive decoders
\end{itemize}

\medskip

\noindent
\textbf{Latent space structure.}
\begin{itemize}
    \item Smoothness enables interpolation
    \item Regularization encourages meaningful representations
\end{itemize}

\medskip

\noindent
\textbf{Evaluation.}
\begin{itemize}
    \item ELBO as training objective
    \item Sample quality and diversity
\end{itemize}

\medskip

\noindent
\textbf{Insight.}  
Practical performance depends on balancing reconstruction, regularization, and model capacity.

\section{A Unifying Perspective}

Variational autoencoders combine several key ideas:

\begin{itemize}
    \item \textbf{Latent variable modeling}
    \item \textbf{Variational inference}
    \item \textbf{Neural network parameterization}
\end{itemize}

\medskip

\noindent
\textbf{Core components.}
\begin{itemize}
    \item Encoder $q_\phi(z \mid x)$
    \item Decoder $p_\theta(x \mid z)$
    \item ELBO objective
\end{itemize}

\medskip

\noindent
\textbf{Key insight.}  
VAEs provide a scalable method for learning probabilistic generative models with latent structure.

\section{Outlook}

This chapter introduced variational autoencoders:
\begin{itemize}
    \item latent representations,
    \item ELBO optimization,
    \item encoder--decoder architectures,
    \item practical considerations.
\end{itemize}

\medskip

In the next chapter, we study generative adversarial networks, which approach generative modeling from a different perspective.

\medskip

\noindent
\textbf{Guiding principle.}  
Combining probabilistic modeling with neural networks enables flexible and scalable generative models.